\scanpage{14}
Nécessité triviale, prouvons la suffisance, i.e. si $M\handwritten{\leadsto}\operatorname{Hom}(M,\mathbb H)$ est
e\deleted{\unsure{d}}xact sur la sous-caté\handwritten{g}orie $\underline C$ de la catégorie des $A$-modules, alors ce
foncteur est exact partout. On sait (Tohoku) qu'il suffit de \deleted{le} vérifier
que pour tout sous-module $N$ de $M=A$, $\operatorname{Hom}(M,H)\handwritten{\to}\operatorname{Hom}(N,H)$ est surjectif,
donc on est ramené au cas où $M,N$ sont \deleted{surjectifs}\inserted{\handwritten{de type fini}}. Or un homom
\[
u:\mathbb N\handwritten{\to}H
\]
a comme image un sous-module de type fini de $H$, donc annulé par une
puissance de $I$, i.e. \inserted{\handwritten{$u$ est}} nul sur un sous-module $I^n\mathbb N$ de $\mathbb N$. En vertu de \uline{Krull},
$I^n\mathbb N$ contient une intersection $N\handwritten{\cap}I^mN$. Or on a une inclusion
\[
0\handwritten{\to}N/N\handwritten{\cap}I^mM\handwritten{\longrightarrow}M/I^mM
\]
\uline{dans} $\underline C$ , et $u$ se factorise par un homomorphisme
\[
u':N/N\; I^mM\handwritten{\longrightarrow}H\quad;
\]
ce dernier par hypothèse peut se prolonger à $M/I^mM$, d'où un homomorphisme
$v:M\handwritten{\to}H$ qui prolonge $u$, cqfd.

\uline{Corollaire} 2.16. Soit $\deleted{\unsure{M}}K$ un $A$-module, et soit $H_I^{\unsure{o|0}}(\deleted{M}K)$ le sous-module de $\deleted{M}K$
formé des éléments annulés par une puissance de $I$. Si $K$ est injectif, il
en est de même de $H_I^o(K)$.

En effet, lorsque $M$ parcourt $\underline C$, on a un isomorphisme fonctoriel
\[
\operatorname{Hom}(M,H_I^o(K))\handwritten{\overset{\sim}{\longrightarrow}}\operatorname{Hom}(M,K)\quad,
\]
de sorte qu'il suffit d'appliquer le critère 2.15.

\deleted{Si on rejette tout ce numéro (comme il me semble raisonnable)}
\uline{Remarques} 2.17. On peut généraliser 2.6. et 2.9. , ainsi que leurs
corollaires, au cas où on est sur un préschéma noethérien $X$, $\underline C$ désignant
la catégorie des Modules cohérents sur $X$ à support $\handwritten{\subset}Y$ ($Y=X$ dans le
cas \deleted{\inserted{dual}} de 2.6.), dans le cas d'un foncteur contravariant \inserted{exact à gauche} $T$ de $\underline C$ dans $(\mathrm{Ab})$:
un tel foncteur \deleted{se représente de façon} se représente (de façon essentiellement unique) comme $\handwritten{F\leadsto}\operatorname{Hom}(F,H)$, où $H$ est un Module quasi-cohérent sur
$X$ à support $\handwritten{\subset}Y$, \margin{\handwritten{déployer}}%
% slanted
encore donné par $\handwritten{H\simeq}\underset{\handwritten{\longrightarrow}}{\lim}_{\handwritten{n}}T(\underline O_X/\underline I^{\handwritten{\unsure{n}}})$. Pour le voir, on note
que $T$ est ``ind-représentable'' par le critère général donné dans A.Grothe\unsure{n}dieck, Technique de descente et théorèmes d'existence en Géom Alg II,
Sém Bourbaki ...., \handwritten{(}qui s'applique ici grâce au fait que les objets de
$\underline C$ sont noethériens\handwritten{),} et on note que $\handwritten{I}\mathrm{nd}\,(\underline C)$ est équivalente de façon naturelle à la catégorie de tous les Modules quasi-cohérents sur $X$ ... Nous
n'aurons sans doute pas besoin de ce fait dans ce Chapitre, où les résultats affines développés ici suffiront pour notre propos.
% archive sheet number: 13

